% ---
% id: reductions/thm:chain
% kind: thm
% title: "Composition of a chain"
% summary: >-
%   The receiver composite of a chain is again a squashing relation with the
%   composed pull-backs, so output-state equality, security transfer and the
%   licence apply to it at every length of the source segment. Its status is
%   exact whenever every link's is.
% prominence: standard
% status: implemented
% uses: [def:chain, def:link, thm:factorization, thm:transfer, cor:sequential, thm:licence, prop:preorder, cor:transported]
% mentions: [rem:ordering, thm:two-sided, req:mixed-status-chain, req:no-approximate-relations]
% assumes: []
% refs: []
% provenance: newly-derived
% implements: certiqs_key.reductions.links:Chain
% lean:
% ---

The order of the source and receiver stages, which this theorem does not
read, is the subject of Remark~\ref{rem:ordering}.

\begin{theorem}\label{thm:chain}
Let $\redChain=(\redChain_{\mathrm{src}},\redChain_{\mathrm{rec}})$ be a
reduction chain with receiver segment
$\redChain_{\mathrm{rec}}=(\Lnk_1,\dots,\Lnk_k)$ and receiver composite
$\Lambda_{\redChain}$ of \eqref{eq:chaincomposite}. Then:
\begin{enumerate}[label=(\roman*),leftmargin=3.0em]
\item $\big(M^{\mathrm{in}}_1;O^{\mathrm{in}}_1\big)
  \dsq_{\Lambda_{\redChain}}
  \big(M^{\mathrm{out}}_k;\hat O^{\mathrm{out}}_k\big)$, and
  $\Lambda_{\redChain}^{\dagger}
  =\Lambda_{\Lnk_1}^{\dagger}\circ\dots\circ\Lambda_{\Lnk_k}^{\dagger}$;
\item Theorem~\ref{thm:factorization}, Theorem~\ref{thm:transfer} and
  Corollary~\ref{cor:sequential} apply to the composite: the physical protocol
  with $M^{\mathrm{in}}_1$ and the final \tref{squashed protocol} with
  $M^{\mathrm{out}}_k$ have the same output state, round by round and
  adaptively;
\item an entry $(\hat O_k,\bowtie_k,c_k)$ declared on the \emph{final} target
  is sound for the receiver composite when the certificate
  \eqref{eq:certificate} holds for its pull-back
  $\Lambda_{\redChain}^{\dagger}(\hat O_k)$, which (C2) computes link by link; and
  Theorem~\ref{thm:licence} then applies with $\Lambda_{\redChain}$;
\item the composite status is \textsf{exact} whenever every link's is, and no
  other composite value is asserted.
\end{enumerate}
Nothing above reads the source segment: (i)--(iv) are statements about the
receiver composite $\Lambda_{\redChain}$ and hold at every length
$L_{\mathrm s}\ge0$. What the \emph{pair} \eqref{eq:chaincomposite} composes to
as a reduction is Theorem~\ref{thm:two-sided}, under the source segment's own
interface conditions and the admission rule of Definition~\ref{def:chain}, and
this theorem does not restate it. Where the source segment is empty the source
composite is the identity and the receiver composite is the whole of the
chain's reduction.
\end{theorem}

\begin{proof}
(i) is Proposition~\ref{prop:preorder} applied $k-1$ times, the tuple chaining
(C2) being exactly the hypothesis that makes the operator tuples compose. (ii)
is those results applied to the relation of (i), which is all they need. (iii)
is Corollary~\ref{cor:transported} and Theorem~\ref{thm:licence} applied to the
composite relation; the pull-back is computed as
$\Lambda_{\Lnk_1}^{\dagger}\big(\dots\Lambda_{\Lnk_k}^{\dagger}
(\hat O_k)\dots\big)$ by (i). (iv): a composition of exact identities is an
exact identity, by (i) applied to identities carrying no error term. The
remaining half is not proved and is not claimed: a chain containing a link
that is not \textsf{exact} is declined rather than assigned a composite
status, by the rules of Requirements~\ref{req:no-approximate-relations}
and~\ref{req:mixed-status-chain}.
The closing paragraph is Definition~\ref{def:chain}: the composite is defined
there as the pair \eqref{eq:chaincomposite}, whose second member is the
receiver segment's and is the only member any of (i)--(iv) reads --- each of
them quantifies over the links $\Lnk_1,\dots,\Lnk_k$ and over nothing else ---
and with $L_{\mathrm s}=0$ its first member is the identity.
\end{proof}
